%% =====================================================================
%%  T-10-06  The Spectacle
%%  -------------------------------------------------------------------
%%  One idea per file. Core is mandatory. Slots in canonical order.
%%  Max 700 words. No layout commands. No filler. New source material
%%  for this entry goes in sources/B3/T-10-06.tex -- never by
%%  rewriting this file (docs/RULES.md R0.1).
%%  =====================================================================
%% @kind: topic
%% @id: T-10-06
%% @title: The Spectacle
%% @subtitle: Striking imagery does the persuading that argument cannot
%% @chapter: c10
%% @order: 60
%% @status: stable
%% @sources: B3
%% @updated: 2026-09-19

\topic{T-10-06}{The Spectacle}{Striking imagery does the persuading that argument cannot}

\begin{core}
Argument reaches the small rational committee in a person; imagery reaches everyone else, immediately. Striking visuals and grand symbolic gestures create an aura of power --- everyone responds to them --- and an audience dazzled by appearances does not audit what is actually being done. The spectacle is attention (T-10-04) with a body: a scene staged so the association you want is felt rather than argued.
\end{core}
\begin{mechanism}
B3's law: stage the scene --- arresting visuals, radiant symbols --- and let the heightened presence do the arguing. The canonical case is Cleopatra's barge --- purple sails, silver oars, gold canopy --- a political negotiation arriving as a sensory event, remembered for two thousand years. The mechanism is pre-attentive processing: symbols land before evaluation begins, and the feelings they install are then mistaken for conclusions the audience reached itself. A spectacle also transfers attributes: associate yourself with one radiating image --- the banner, the building, the signature look --- and the image's qualities (scale, permanence, nerve) are credited to you by the audience that never checks the transfer. The reversal in the source is total: power never flows from ignoring images and symbols, so none is offered.
\end{mechanism}
\begin{conditions}
Strongest at moments of first judgment --- pitches, launches, entrances, trials --- and for claims that cannot be demonstrated directly (competence, stability, inevitability). It weakens with repeated exposure (the same spectacle decays into scenery), among expert audiences who read staging as a cost signal to be audited, and anywhere the gesture outruns the substance it promises --- the audience eventually meets the substance, and the meeting is the invoice.
\end{conditions}
\begin{application}
  \item Choose one symbol and keep it. The spectacle that repeats becomes identity; scattered gestures become noise.
  \item Stage for the claim you cannot argue: scale for stability, ritual for legitimacy, one outrageous visual for nerve.
  \item Spend the audience's sensory memory --- one scene, vivid, retellable. If they cannot retell it, it did not happen.
  \item Match the register of the room: spectacle that over-reaches the occasion reads as desperation wearing costume.
  \item Have the substance ready to meet the converts. A spectacle is a bridge; traffic must find ground on the far side.
\end{application}
\begin{feedback}
  \item Working: your event is retold by people who were not there --- with the numbers wrong but the feeling right.
  \item Working: the association appears in how strangers introduce you.
  \item Failing: compliments arrive for the production, not the claim. The scene ate the message.
  \item Failing: the spectacle is mocked for trying. Register missed; the symbol now works against the substance.
\end{feedback}
\begin{failure}
Spectacle is expensive in both directions: it costs real resources to stage and real credibility when it flops. The deeper failure is substitution --- operators who can stage begin to prefer staging, and the gap between image and substance widens until the audience's eventual audit closes it destructively. And symbols are hostage to their contexts: one scandal re-prices the banner you built your presence on.
\end{failure}
\begin{countermeasures}
  \item When moved by a presentation, name the symbols to yourself. The spell weakens under labelling --- and you learn the operator by their prop list.
  \item Audit image-to-substance ratios in your own house: what does the flagship event cost, and what does it change? The gap is theatre, and theatre has its place --- named.
  \item Watch for borrowed symbols: the flag, the lab coat, the corner office rented for the day (T-18-01). Symbols are cheap; check what they are attached to.
  \item Retell-resistant fact-check: whatever survives retelling intact is substance. What grows in retelling is spectacle.
\end{countermeasures}

\sources{B3}
\seesources
\seealso{T-10-04, T-17-05, T-18-01}
